\subsection{从模型图到运行时分配}
模型准备和 Provider 分配发生在不同时刻。准备阶段识别支持的模型、张量接口及候选切分，并可在具体调用之前生成描述分片和依赖的可复用模板。运行时规划再根据请求和收集的 positive ACK 将角色绑定到实际 Provider。手工模板符合这种设计；检验协议并不要求自动搜索切分。Controller 的授权职责也不使它成为模型切分器。

\begin{figure}[htbp]\singlespacing\centering
\begin{tikzpicture}[every node/.style={draw,align=center,font=\small,text width=9.6cm,inner sep=6pt}]
\node (model) at (0,0) {模型图及支持的张量接口};
\node (template) at (0,-1.25) {准备好的分片及依赖模板};
\node (assign) at (0,-2.85) {请求要求 + 通过验证的 positive ACK\newline 运行时 Provider 分配};
\node (plan) at (0,-4.55) {已提交角色及命名依赖\newline Provider 本地输入检查};
\node (execute) at (0,-6.25) {准备被分配分片、交换获准张量\newline 执行并发布结果};
\draw[->] (model)--(template);\draw[->] (template)--(assign);
\draw[->] (assign)--(plan);\draw[->] (plan)--(execute);
\end{tikzpicture}
\caption{模型准备与请求专属协作的拟议分工。箭头表示处理顺序，并不声称所有模型图和分配策略都已实现。}
\end{figure}

各分片契约需要定义输入输出张量名称、形状或支持的动态形状约束、元素类型、序列化及模型版本兼容性。跨越跳跃连接或分支图的切分可能同时需要多个输出，而非单个向量。接收分片必须得到全部必要输入才能执行。一个有效反例是提交签名正确、但形状或模型身份错误的张量：协议来源有效不意味着应用契约成立。

Pipeline partitioning 分配连续计算阶段并传递边界张量。Tensor parallelism 将操作分配给多个参与者，并增加 reduction、gather 等交换。模型复制则由完整副本承担独立工作。这些是不同执行策略，不是把模型放到多台设备的同义词。Early exit 改变计算量，也应与放置分别评价。ONNX 导出提供图表示~\cite{onnx}，可用切分及 collective 仍取决于 runtime 和模型算子的支持。

当 Provider 缺少所需对象或兼容的已加载会话时，准备代价尤为重要。评价应区分对象传输、会话创建、加速器内存分配和执行。复用驻留权重的 warm run 不能在不披露差异时与 cold baseline 比较。放置策略可以偏好已持有兼容状态的 Provider，但必须遵守请求授权，不能假定缓存状态可以用于任意新请求。

\subsection{张量交换与自回归状态}
在流水线中，上游分片生成激活张量，使其作为受保护命名 Data 可被获取，并通过声明依赖提供引用。下游取得和验证对象、检查张量契约，再执行本分片。Fan-out 向多个获准消费者复用上游输出；fan-in 等待计算所需输入集合。协作协议判断哪些生产者和依赖边获准，DI 判断数值和布局是否适合模型。

自回归推理增加了解码步骤之间的依赖。启用 K/V caching 时，prefill 后的计算使用前一 token 与保留的 K/V 状态，避免重新计算整个前缀。按层切分时，各 runtime 可以保留本分片层的 K/V 张量。生成的 token 或选择它所需的信息返回驱动下一解码步骤的组件。因此，一次前向计算可以沿有向无环图流动，而完整生成过程会随时间反复执行该图。

权重驻留加速器并不消除逻辑上的 K/V 输入。Runtime 可以复用设备缓冲区并直接绑定，避免每步序列化和传输整个缓存~\cite{ortkvcache}。重新分配或会话继续需要明确的兼容性及状态转移策略；以无关 K/V 状态启动替换者不是有效恢复。评价应区分本地 K/V 复用收益、远程转移成本和框架编排开销。

初始数值验证应将支持的切分执行与相同模型、精度、输入和解码策略下的未切分执行比较，并在测量前规定张量容差或确定性输出判据。小型合成模型能够验证张量路径，但不能证明大型预训练模型的内存可行性、吞吐量或恢复。这些控制区分通信路径可用与分布式推理正确且有价值。
